\documentclass[12pt, landscape, oneside]{book}

% --- Paquetes Críticos ---
\usepackage[utf8]{inputenc}
\usepackage[T1]{fontenc}
\usepackage[spanish]{babel}
\usepackage{geometry}
\geometry{
    a4paper,
    landscape,
    top=2cm,
    bottom=2cm,
    left=2cm,
    right=2cm,
    headheight=15pt
}

\usepackage{xcolor}
\usepackage{tikz}
\usetikzlibrary{shapes.geometric, arrows, positioning, shadow, patterns, calc, trees}
\usepackage{tcolorbox}
\tcbuselibrary{skins, breakable, listings}
\usepackage{titlesec}
\usepackage{fancyhdr}
\usepackage{graphicx}
\usepackage{hyperref}
\usepackage{amsmath}
\usepackage{amssymb}
\usepackage{booktabs}
\usepackage{multicol}
\usepackage{enumitem}

% --- Configuración de Colores Corporativos ---
\definecolor{CorporateBlue}{RGB}{21, 101, 192} 
\definecolor{CorporateGold}{RGB}{255, 160, 0} 
\definecolor{SoftGray}{RGB}{245, 245, 245}
\definecolor{DeepDark}{RGB}{33, 33, 33}

% --- Estilo de Títulos ---
\titleformat{\chapter}[display]
  {\normalfont\huge\bfseries\color{CorporateBlue}}
  {\chaptertitlename\ \thechapter}{10pt}{\Huge}
\titlespacing*{\chapter}{0pt}{-20pt}{30pt}

% --- Cajas de Usuario (Friendly) ---
\newtcolorbox{pasos}[1]{
  colback=white,
  colframe=CorporateBlue,
  fonttitle=\bfseries,
  title=Pasos para: #1,
  enhanced,
  attach boxed title to top left={xshift=10pt, yshift=-10pt},
  boxed title style={colback=CorporateBlue}
}

\newtcolorbox{escenario}[1]{
  colback=CorporateGold!5,
  colframe=CorporateGold,
  fonttitle=\bfseries\color{black},
  title=Caso Real: #1,
  enhanced,
  boxed title style={colback=CorporateGold}
}

% --- Documento ---
\begin{document}

% --- Portada ---
\begin{titlepage}
    \centering
    \vspace*{3cm}
    {\fontsize{50}{60}\selectfont\bfseries\color{CorporateBlue} GANESHA ERP}\\[1cm]
    {\huge\bfseries\color{CorporateGold} Manual de Usuario: Guía de Referencia Total}\\[2cm]
    
    \begin{tikzpicture}
        \draw[CorporateBlue, ultra thick, rounded corners] (0,0) rectangle (15,4);
        \node at (7.5,3) {\Large\bfseries "Diseñado para ser entendido por humanos,"};
        \node at (7.5,2) {\Large\bfseries "ejecutado por profesionales."};
    \end{tikzpicture}
    
    \vfill
    {\large FlasheyEstudi | 2026}
\end{titlepage}

\tableofcontents

% ==========================================
\chapter{Primeros 15 Minutos: Inicio Rápido}
% ==========================================

Si es tu primera vez en Ganesha, sigue este camino para tener el sistema listo en minutos:

\begin{pasos}{Configuración Relámpago}
    \begin{enumerate}
        \item \textbf{Crea tu Empresa:} En el menú lateral, ve a \textbf{Empresas}. Pon tu RUC, Nombre y Moneda.
        \item \textbf{Verifica tu Catálogo:} Ve a \textbf{Cuentas}. El sistema ya trae un catálogo base. Ajusta solo lo necesario.
        \item \textbf{Abre el Mes:} En \textbf{Períodos}, asegúrate de que el mes actual esté en estado "ABIERTO".
        \item \textbf{¡Listo!:} Ya puedes registrar tu primera factura o movimiento.
    \end{enumerate}
\end{pasos}

% ==========================================
\chapter{Diccionario: Del "Sistemas" al "Español"}
% ==========================================

Para que nos entendamos, aquí tienes qué significa cada cosa técnica que podrías ver:
\begin{description}
    \item[Dashboard:] Tu tablero de control con gráficas.
    \item[Third Party (Tercero):] Cualquier persona de afuera (Cliente o Proveedor).
    \item[Journal Entry (Póliza):] Un registro contable (quién le dio a quién).
    \item[Audit Log:] El historial de quién tocó qué.
    \item[Asset (Activo):] Algo que compraste y que dura mucho (una laptop, un carro).
\end{description}

% ==========================================
\chapter{Explicación Detallada de los Módulos}
% ==========================================

\section{Módulo: Facturación (Ventas y Compras)}
\textbf{¿Para qué sirve?}: Para anotar lo que vendes y lo que compras.
\begin{itemize}
    \item \textbf{Acción}: Creas una factura a un cliente.
    \item \textbf{Repercusión Real}: El sistema le "avisa" al módulo de \textbf{Cuentas por Cobrar} que ese cliente te debe dinero. También le avisa al \textbf{Dashboard} para que suba tu gráfica de ingresos.
\end{itemize}

\section{Módulo: Bancos y Caja}
\textbf{¿Para qué sirve?}: Para que tu saldo en el sistema coincida con tu estado de cuenta real.
\begin{itemize}
    \item \textbf{Acción}: Registras que el cliente te pagó la factura.
    \item \textbf{Repercusión Real}: El dinero entra a tu cuenta de banco en el sistema y "limpia" la deuda del cliente en el módulo de Facturación.
\end{itemize}

\section{Módulo: Activos Fijos (Tus Bienes)}
\textbf{¿Para qué sirve?}: Para que no pierdas de vista cuánto valen tus cosas con el tiempo.
\begin{itemize}
    \item \textbf{Acción}: Registras una nueva camioneta.
    \item \textbf{Repercusión Real}: Cada mes, el sistema calculará cuánto valor perdió (depreciación) y lo enviará al módulo de \textbf{Gastos} automáticamente.
\end{itemize}

% ==========================================
\chapter{Casos de la Vida Real (Survival Guide)}
% ==========================================

\begin{escenario}{"Compré una computadora para la oficina"}
    \begin{enumerate}
        \item Ve al módulo de \textbf{Facturación} y registra la compra al proveedor.
        \item Al guardar, el sistema genera la deuda al proveedor.
        \item Ve al módulo de \textbf{Activos Fijos} y registra la computadora para que empiece a depreciarse cada mes.
        \item Cuando le pagues al proveedor, usa el módulo de \textbf{Bancos} para "matar" esa deuda.
    \end{enumerate}
\end{escenario}

\begin{escenario}{"Quiero saber si estoy ganando dinero"}
    \begin{enumerate}
        \item Mira tu \textbf{Dashboard}. Ahí ves la gráfica de Ingresos vs Egresos.
        \item Si quieres el detalle oficial, ve al módulo de \textbf{Reportes} y saca un "Estado de Resultados".
        \item Si tienes dudas de un número, pregúntale a la \textbf{IA (Ganesha AI)}: "¿Por qué gasté tanto en papelería este mes?".
    \end{enumerate}
\end{escenario}

% ==========================================
\chapter{Matriz de Interacción: El Camino del Dato}
% ==========================================

Para los que quieren entender el "cableado" interno:

\begin{center}
\begin{tikzpicture}[node distance=1.5cm]
    \node (entry) [draw, fill=blue!10, rounded corners] {\textbf{Tú ingresas un dato (Factura/Banco)}};
    \node (logic) [draw, fill=orange!10, below=of entry] {\textbf{El Sistema aplica reglas (Impuestos/IVA)}};
    \node (cont) [draw, fill=green!10, below=of logic] {\textbf{Se crea una Póliza Automática (Contabilidad)}};
    \node (rep) [draw, fill=purple!10, below=of cont] {\textbf{Tus Reportes y Gráficas se actualizan}};
    
    \draw[->, thick] (entry) -- (logic);
    \draw[->, thick] (logic) -- (cont);
    \draw[->, thick] (cont) -- (rep);
\end{tikzpicture}
\end{center}

\vfill
\begin{center}
    \color{gray} \textit{Este manual fue diseñado para que tú te enfoques en tu negocio, mientras Ganesha se encarga de los números.}
\end{center}

\end{document}
